\section{Empirical results}
\subsection{Completion rates and admissible observations}
Of 2813 scheduled computational jobs, 57 terminated without a successful numerical record. The resulting dataset comprises 92,679 measured or explicitly classified method-level outcomes. A total of 76,220 (82.24\%) satisfy the predefined completeness and status conditions. The remaining 16,459 comprise 11,208 numerical-rejection records, 4224 additional smoothness-inversion failure records, 970 computational-budget exclusions, and 57 records representing unavailable job inputs. The classification retains unsuccessful searches; omitting them would bias subsequent comparisons toward implementations with a higher apparent completion rate.

Figure~\ref{fig:reliability} displays stacked status fractions by criterion and implementation, with distinct colors identifying the methods and patterned segments indicating numerical rejection, inversion failure, or resource exclusion. Table~\ref{tab:empirical-reliability} reports the underlying counts by method. Because eligibility varies across configurations, these percentages summarize observed method-level records rather than randomized or matched failure probabilities.

\begin{figure}[t]
\centering
\includegraphics[width=\linewidth]{\VizDir/figures/fig01_reliability_stacked.pdf}
\caption{Stacked status composition of complete-criterion evaluations by selection criterion. Solid colors identify the implementation, while patterned segments distinguish the recorded failure mechanisms. Rates are conditional on method-specific eligible observations and do not constitute matched failure probabilities.}
\label{fig:reliability}
\end{figure}

\input{\VizDir/tables/table02_reliability.tex}

\subsection{Complete-grid running time}
Table~\ref{tab:empirical-paired} summarizes the matched cold-time ratios for eight complete 100-candidate CV and GCV searches, reporting both seed-level median and interquartile range and the previously calculated geometric mean with seed-cluster bootstrap intervals. Figure~\ref{fig:paired} displays the corresponding distributions as violin densities, narrow boxplots and individual seed-level paired ratios. Each case uses identical observations and candidates for both methods. Ratios $\rho=T_{\mathrm B}/T_{\mathrm S}>1$ favor spectral SVD; all eight representative grids show this advantage. The distributions reveal variation and asymmetry hidden by the point summaries and must not be interpreted as eight identically distributed experiments.

\input{\VizDir/tables/table01_runtime.tex}

\begin{figure}[t]
\centering
\includegraphics[width=\linewidth]{\VizDir/figures/fig02_runtime_distributions.pdf}
\caption{Observed distributions of matched banded-Cholesky/SVD complete cold-time ratios over 100 smoothing candidates. The violin outlines are descriptive kernel-density estimates, the boxes give medians and interquartile ranges, and the dots represent independent-seed cluster summaries. The dashed line denotes equal time; numerical confidence intervals for geometric means are given in Table~\ref{tab:empirical-paired}.}
\label{fig:paired}
\end{figure}

The effect of candidate count is particularly clear for conditional CV at $n=256$ and $d=3$. The banded/SVD time ratio is 0.698 (95\% interval 0.671--0.724) for one candidate and one series, but 4.846 (4.567--5.125) for 100 candidates. For 100 candidates evaluated on three series, the corresponding ratio is 3.655 (3.450--3.860). The reversal is consistent with a sizable spectral setup cost that is recovered through repeated criterion evaluation. It does not imply a universal crossover at $M=100$; both the location and even the existence of a crossover depend on the criterion, series length, order and implementation.

For fixed-mean first-order GCV at $n=64$, the analogous ratio is 1.001 (0.982--1.018) for one candidate and 2.328 (2.282--2.374) for the 100-candidate grid. Thus a negligible single-candidate performance difference can coexist with a substantial grid-search advantage.

\subsection{Reuse across smoothing coordinates and series}
The candidate representation also influences total cost. Two computational policies were compared: repeated conversion of smoothness values $S$ to the penalty $\lambda$, and a single conversion shared across the search. The candidate set is identical under both policies, so the contrast concerns computational reuse rather than a different optimization problem.

For spectral CV with $n=256$, $d=3$, 100 candidates and three series, the geometric complete cold-time ratio of once-mapped penalties to repeated direct-$S$ conversion is 0.580 (95\% interval 0.570--0.590). In this setting, precomputing the candidate transformation reduces total execution time. Figure~\ref{fig:mapping} displays corresponding descriptive comparisons across other numerically admissible multi-series grid configurations.

\begin{figure}[t]
\centering
\includegraphics[width=\linewidth]{\VizDir/figures/fig03_mapping_distributions.pdf}
\caption{Per-case-seed distributions of mapped-penalty/direct-smoothness cold-time ratios, grouped by criterion, method and candidate-grid size. Violin density and boxplots describe heterogeneous admissible cases under the common execution phase. Values below one favor mapping the penalty grid once; the pooled descriptive distributions are not universal algorithmic scaling laws.}
\label{fig:mapping}
\end{figure}

Figure~\ref{fig:workloads} illustrates how a fixed first-order GCV configuration responds to four combinations of candidate and series count. The comparisons show why initialization, single-candidate scoring, multiple-series reuse and full-grid selection must be distinguished when assessing an implementation. In particular, a procedure optimized only for one fitted trend need not minimize the cost of selecting a smoothing parameter.

\begin{figure}[t]
\centering
\includegraphics[width=\linewidth]{\VizDir/figures/fig04_workload_stacked.pdf}
\caption{Stacked component-wise medians for initialization, smoothness-to-penalty mapping and criterion evaluation in a first-order GCV example ($n=64$, $M=100$). Colors identify the implementation; light, medium and dark sections identify stages. Workload labels denote number of smoothing candidates times number of series. Sums of component medians are not generally the medians of total cold time.}
\label{fig:workloads}
\end{figure}

\subsection{Numerical reliability and its consequences}
The aggregate status proportions in the 24-worker execution phase are 10,848/13,176 (82.3\%) for banded CV and 13,176/13,176 (100\%) for spectral CV. The corresponding GCV counts are 13,176/16,280 (80.9\%) and 15,892/16,086 (98.8\%), respectively. Their denominators differ because method eligibility and the evaluated configurations differ; they should not be used as an unconditional estimate of solver accuracy. The comparison nevertheless motivates a closer examination of the failure mechanisms, reported next.

\subsection{Post-solve optimality failures in banded computation}
The banded-Cholesky implementation recorded 9096 method-level
timing observations rejected by the scaled stationarity-residual
check in Eq.~\eqref{eq:kkt}. They represent 3384 observations
under CV, 4160 under GCV, and 1552 under marginal likelihood.
Among the \emph{rejected} records, the smallest reported value
of $\eta_{\mathrm{opt}}$ was $1.0013\times10^{-7}$, the median was
$1.4079\times10^{-7}$, and the largest was
$1.0587\times10^{-6}$, against the specified threshold
$10^{-7}$. Figure~\ref{fig:kktresidualdist} displays this
conditional distribution. The data do not record the
corresponding residual magnitudes for accepted runs; the
histogram must not be used as an estimate of the full
residual distribution or of an independent failure probability.
Since the check follows factorization and solution, the
failures establish rejection by the \emph{complete banded
criterion pipeline}, not an inability of Cholesky to
factor an SPD matrix.

\begin{figure}[t]
\centering
\includegraphics[width=\linewidth]{\KKTDir/kkt_rejection_distribution.pdf}
\caption{Distribution of the scaled first-order optimality
residual $\eta_{\mathrm{opt}}$ among \emph{rejected} banded
complete-criterion computations. Histogram segments are
colored by statistical criterion (CV, GCV and Gaussian
marginal likelihood). The dashed line marks the
implementation-specific $10^{-7}$ acceptance threshold;
the dotted line marks the median of rejected records.
Counts are correlated method-level timing observations,
not independent failures. Residuals of accepted
evaluations are not available in the archived record.}
\label{fig:kktresidualdist}
\end{figure}

A particularly informative contrast occurs in nine parameter configurations. For each configuration, 97 independent realizations were considered. Both methods completed every recorded single-candidate evaluation, but the banded implementation rejected every complete-grid evaluation, whereas the spectral implementation completed the corresponding grids. The configurations comprise CV and GCV with $(n,d,M)=(128,3,1000)$, $(256,4,100)$ and $(256,3,1000)$; GCV with $(64,3,10000)$; and profiled Gaussian marginal likelihood with $(256,3,100)$ and $(512,3,100)$. Figure~\ref{fig:cholgridfail} displays these results.

The nine configurations account for 6984 banded-rejection timing records; another 2112 such records arise from larger fourth-order configurations. These are repeated method-level observations within 97 independent seed clusters per configuration, rather than thousands of independent experimental realizations. The saved status information does not identify the exact smoothing-grid point that first produced each rejection, and hence no candidate-specific instability boundary is inferred.

\begin{figure}[t]
\centering
\includegraphics[width=\linewidth]{\VizDir/figures/fig05_grid_completion.pdf}
\caption{Full-grid numerical completion in nine parameter configurations, with method-specific coloring and patterned rejection fractions. All recorded single-candidate outcomes complete for both implementations, but the banded full-grid pipeline is rejected while the spectral pipeline completes. Each configuration has 97 independent seed clusters; within-seed repeats, mapping policies and multiple-series workloads are correlated.}
\label{fig:cholgridfail}
\end{figure}

\input{\VizDir/tables/table03_grid_failures.tex}

The apparent spectral advantage must be interpreted as \emph{observed computational completion}, not a controlled comparison of identical backward-error tests. The spectral scorer does not impose the same implementation-specific acceptance criterion as the banded scorer. Independent numerical reference calculations would be required to establish a strict relative-accuracy ranking at each penalty.

\subsection{Smoothness inversion remains a common numerical limitation}
An additional failure mechanism affects both implementations at larger difference order and series length. In the fourth-order CV/GCV configurations with $n=1024$, 2048 and 4096, the numerical inversion of the normalized smoothness mapping sometimes returned values outside its prescribed tolerance. Across the evaluated configurations, 6336 timing records explicitly describe inaccurate smoothness inversion: 2112 associated with the banded method and 4224 with SVD. The distribution of observed outcomes in the fourth-order settings is given in Fig.~\ref{fig:fourthfail}.

This error concerns the transformation from the bounded coordinate $S$ to the penalty $\lambda$, rather than evidence that a spectral decomposition failed. Crucially, a successfully completed one-candidate evaluation cannot be used as a substitute for an incomplete multi-candidate search. Accordingly, the large-order configurations with shared inversion difficulties do not supply an empirical winner for the full criterion grid.

\begin{figure}[t]
\centering
\includegraphics[width=\linewidth]{\VizDir/figures/fig06_order4_failures.pdf}
\caption{Stacked method-colored fractions of numerically complete evaluations and distinct rejection mechanisms for difference order four. The bands show implementation-specific complete, post-solve KKT-rejected and inaccurate smoothness-inversion outcomes. Denominators consist of available method-level records; both implementations exhibit inversion failures at higher sample lengths.}
\label{fig:fourthfail}
\end{figure}

\subsection{The single-candidate versus grid-search crossover}
The principal runtime trade-off becomes visible when paired, numerically admissible observations are compared across candidate counts. Figure~\ref{fig:onegrid} shows three examples evaluated under identical matching rules. In conditional CV with $n=256$, $d=3$, the geometric banded/SVD cold-time ratio changes from 0.698 (95\% interval 0.671--0.724) for one candidate to 4.846 (4.567--5.125) for 100 candidates. In fixed-mean GCV with $n=64$, $d=1$, the corresponding ratios are 1.001 (0.982--1.018) and 2.328 (2.282--2.374). For fixed-mean CV with $n=128$, $d=3$, they are 1.069 (1.052--1.087) and 5.401 (5.244--5.555).

\begin{figure}[t]
\centering
\includegraphics[width=\linewidth]{\VizDir/figures/fig07_crossover_distributions.pdf}
\caption{Distributional comparison of matched banded-Cholesky/SVD cold-time ratios for one candidate and 100 candidates. Violin densities, boxplots and individual seed-level points show the pronounced reversal of computational preference. Values above one favor SVD. Corresponding independent-seed bootstrap intervals for geometric means are reported in the text rather than inferred from violin widths.}
\label{fig:onegrid}
\end{figure}

These results support the computational interpretation of Eqs.~\eqref{eq:spectraltrend}--\eqref{eq:rss}: spectral setup is paid once, while the attenuation coefficients and the statistics needed for subsequent criterion evaluations are reused. The crossover is observed in complete scoring workloads; it should not be conflated with the cost of fitting a trend for a prescribed penalty.

\subsection{Memory requirements and computational feasibility}
The principal disadvantage of the spectral approach is its dense storage requirement. For the evaluated implementation, the conservative memory admission allowance corresponds to five $n\times n$ double-precision arrays, or $40n^2$ bytes. An explicit dense-inverse reference is assigned seven such arrays, or $56n^2$ bytes. In contrast, a banded Cholesky factor and a banded selected-inverse buffer together require approximately $16n(d+1)$ bytes, excluding ancillary data structures. Figure~\ref{fig:memorytradeoff} illustrates these storage relations for $d=4$.

For example, the spectral admission allowance is 640~MiB at $n=4096$. At $n=8192$ it rises to 2.5~GiB, exceeding the 2~GiB dense allocation limit imposed in the study. The latter sample length was not evaluated in the present comparison. The curves therefore represent \emph{implementation-specific memory allowances and illustrative buffer sizes}, not measurements of peak resident memory or empirical allocation failures.

\begin{figure}[t]
\centering
\includegraphics[width=\linewidth]{\VizDir/figures/fig08_memory_scaling.pdf}
\caption{Analytical source-code memory allowances as a function of series length for dense spectral SVD, an explicit inverse and two fourth-order banded working arrays. The shaded region denotes examined sample lengths and the dashed line the 2~GiB dense allocation policy. These curves are storage-policy calculations, not measured peak memory.}
\label{fig:memorytradeoff}
\end{figure}

There were 970 resource-budget exclusions in the recorded timing outcomes: 776 for the explicit-inverse reference and 194 for SVD. These were induced by prescribed computational-work thresholds, not observed out-of-memory errors. Thus the memory analysis establishes a concrete feasibility trade-off without attributing numerical failures to unsupported memory measurements.


\section{Discussion}
The empirical evidence supports a clear implementation preference, subject to mathematical and computational feasibility: \emph{reuse of a spectral decomposition is the recommended default for finite-difference PLS smoothing-parameter selection over a moderate-to-large candidate grid}. In contrast to a sequence of penalized normal-equation factorizations, spectral preprocessing represents the complete search in one orthogonal basis. It amortizes setup, supplies degrees of freedom and residual sums of squares directly, and provides a common representation for multiple selection criteria. The advantage is pronounced in the matched CV/GCV grid comparisons and is compatible with the observed numerical-completion differences.

The recommendation should not be extended beyond the observed evidence. First, the banded method is often preferable when evaluating only one penalty. Second, the spectral basis and its work arrays require quadratic memory, which can dominate the decision for long series even when SVD is attractive at moderate $n$. Third, numerical acceptance is not interchangeable with certified forward accuracy: the two implementations apply different diagnostic checks, and both encounter smoothness-inversion difficulties in the larger fourth-order cases. Fourth, the recorded multi-process execution phases limit the transportability of wall-clock ratios to idle workstations or different parallel workloads.

These limitations suggest the following practical selection principle. If the task is a multi-candidate \emph{complete-criterion} search and the dense spectral workspace satisfies a conservative memory budget, prepare the spectral decomposition once, reuse series-specific projections and precompute the candidate penalties when appropriate. If only one candidate is needed, or if dense allocation would be prohibitive, a structured banded method is a reasonable alternative subject to numerical residual checks. When a requested smoothness value cannot be inverted accurately, neither method should silently substitute an approximate grid winner: the affected comparison must be marked invalid or handled by a separately validated numerical strategy.

A further limitation is the restriction to synthetic data and one computational environment. Seed-cluster intervals quantify variability under the observed experimental protocol, but not variability across operating systems, processors, library versions or CPU utilization. The validity masks can induce survivor bias in timing ratios conditional on completed searches. The study does not establish improvements in trend-estimation risk or forecasting accuracy, and it does not constitute an independent high-precision verification of every evaluated penalty. Future work should isolate the source of the large-order smoothness-inversion failures, independently certify residual and criterion accuracy, and compare implementations under controlled processor occupancy.

\section{Conclusions}
This study establishes that the preferred computational method for penalized least-squares trend smoothing depends on the \emph{parameter-selection workload}, not merely the cost of one linear solve. Across valid matched 100-candidate CV and GCV searches, spectral SVD consistently outperformed banded Cholesky in the representative tested configurations. A measured conditional-CV setting changed from favoring banded Cholesky for one candidate to favoring spectral SVD by a factor of approximately 4.85 for the complete grid. The numerical-status analysis further identified repeated banded post-solve failures on several complete grids, while retaining documented failure modes of the spectral pipeline.

We therefore recommend \emph{spectral decomposition as the primary implementation for complete, repeated smoothing-parameter selection whenever its dense storage is affordable}. Banded methods remain valuable for isolated evaluations and long series, and numerical diagnostics remain mandatory for either implementation. The contribution is a measured, criterion-aware recommendation for computational practice rather than a new trend estimator or an unconditional claim of spectral superiority.

\section{Code and data availability}
The source implementations, analysis programs, selection-status records, derived timing data and reproducible figures are available in the public project repository at \url{https://github.com/heri-espino/trendestimation}. The underlying raw execution records are retained by the author but are not publicly deposited; their availability is subject to author approval. The publicly archived derived observations preserve incomplete and numerically inadmissible searches, and include the pairing and bootstrap metadata required to reproduce the tabulated comparisons.
